\documentclass[11pt,oneside]{memoir}

% Packages
\usepackage{amsmath, amssymb, amsthm}
\usepackage{tikz}
\usetikzlibrary{cd}
\usepackage{hyperref}
\usepackage{geometry}
\geometry{margin=1.2in}

% Theorem Environments
\theoremstyle{definition}
\newtheorem{definition}{Definition}[chapter]

\theoremstyle{plain}
\newtheorem{theorem}[definition]{Theorem}
\newtheorem{lemma}[definition]{Lemma}
\newtheorem{proposition}[definition]{Proposition}

\theoremstyle{remark}
\newtheorem{remark}[definition]{Remark}

% Title
\title{Anima-Vectorium v2.0 Mathématique \\ 
\large A Formal Expressive--Geometric Architecture}
\author{Borealis S. Hedling}
\date{2026}

\begin{document}
\maketitle

\tableofcontents

\chapter{Introduction}

Anima-Vectorium is a governed expressive--geometric architecture formalized as a 
smooth manifold system equipped with operator algebra, deterministic kernel flow, 
and a homotopy-invariant safety region. Version 2.0 Mathématique provides the 
full mathematical formalization of the system, including manifold definitions, 
operator algebra, curvature tensors, holonomy groups, category-theoretic adjunctions, 
and verification invariants.

\chapter{Expressive Manifold}

\section{Definition}

\begin{definition}[Expressive Manifold]
Let $E_M$ be a smooth, second-countable, Hausdorff manifold with atlas 
$\mathcal{A}_E$. An expressive chart is a smooth map 
$u : U \subset E_M \to \mathbb{R}^n$ such that transitions 
$u \circ v^{-1}$ are $C^\infty$ and expressive fibers form a trivial bundle.
\end{definition}

\section{Tangent Structure}

The expressive tangent bundle $T(E_M)$ is defined in the usual manner. Expressive 
vectors represent directional expressive change.

\chapter{Projection Functor}

\section{Functorial Structure}

\begin{definition}[Projection Functor]
Define categories $\mathbf{Expr}$ and $\mathbf{Geom}$. A functor 
$\Pi : \mathbf{Expr} \to \mathbf{Geom}$ preserves composition and identity:
\[
\Pi(f \circ g) = \Pi(f) \circ \Pi(g), \qquad \Pi(\mathrm{id}_E) = \mathrm{id}_{\Pi(E)}.
\]
\end{definition}

\section{Naturality}

\begin{theorem}[Naturality of Projection]
For any expressive morphism $f : E_1 \to E_2$,
\[
\Pi(f) \circ \eta_{E_1} = \eta_{E_2} \circ f,
\]
where $\eta$ is the expressive--geometric natural transformation.
\end{theorem}

\chapter{Geometric Manifold}

\section{Structure}

\begin{definition}[Geometric Manifold]
Let $G_M$ be a smooth Riemannian manifold with metric $g$ and Levi-Civita 
connection $\nabla$.
\end{definition}

\section{Operator Domains}

\begin{definition}[Operator Domain]
An operator $O : T(G_M) \to T(G_M)$ is smooth and linear.
\end{definition}

\chapter{Operator Algebra}

\section{Lie Algebra}

Define operators:
\[
\mathcal{O} = \{\mathrm{Hol}, \mathrm{Curv}, \mathrm{Drift}, \mathrm{Stab}\}.
\]

\begin{definition}[Operator Lie Algebra]
The Lie algebra $\mathfrak{A}$ is generated by $\mathcal{O}$ with bracket 
$[X,Y] = XY - YX$.
\end{definition}

\section{Commutation Relations}

\begin{theorem}
$[\mathrm{Hol}, \mathrm{Curv}] \neq 0$.
\end{theorem}

\begin{theorem}
$[\mathrm{Drift}, \mathrm{Stab}] = 0$.
\end{theorem}

\section{Operator Algebra Table}

\begin{center}
\begin{tabular}{|c|c|c|}
\hline
Operator & Domain & Notes \\
\hline
Hol & $T(G_M)$ & Non-commutative with Curv \\
Curv & $T(G_M)$ & Tensorial curvature operator \\
Drift & $T(G_M)$ & Commutes with Stab \\
Stab & $T(G_M)$ & Stability-preserving \\
\hline
\end{tabular}
\end{center}

\chapter{Kernel Flow}

\section{Deterministic Flow}

\begin{definition}[Kernel Flow]
Let $K : G_M \to G_M$ be generated by vector field $X$:
\[
\frac{d}{dt} \phi_t(g) = X(\phi_t(g)).
\]
\end{definition}

\begin{theorem}[Determinism]
If $\phi_t(g_1) = \phi_t(g_2)$ for all $t$, then $g_1 = g_2$.
\end{theorem}

\chapter{Safety Homotopy}

\section{Safety Region}

\begin{definition}[Safety Region]
Let $\mathrm{SafeRegion} \subset G_M$ be a closed embedded submanifold.
\end{definition}

\section{Homotopy}

\begin{definition}[Safety Homotopy]
A homotopy $H : G_M \times [0,1] \to G_M$ is safety-preserving if 
$H(g,t) \in \mathrm{SafeRegion}$.
\end{definition}

\begin{theorem}[Homotopy Invariance]
If $g_0 \sim g_1$ via a safety homotopy, then both lie in $\mathrm{SafeRegion}$.
\end{theorem}

\section{Safety Diagram}

\begin{center}
\begin{tikzpicture}
\draw (0,0) ellipse (3 and 2);
\node at (0,0) {SafeRegion};
\draw[->] (-4,0) -- (-1,0);
\draw[->] (1,0) -- (4,0);
\node at (-4.5,0) {$g_0$};
\node at (4.5,0) {$g_1$};
\end{tikzpicture}
\end{center}

\chapter{Category-Theoretic Structure}

\section{Adjunction}

\begin{theorem}[Adjunction]
\[
\Pi \dashv K
\]
with natural bijection:
\[
\mathrm{Hom}_{\mathbf{Geom}}(\Pi(E), G) 
\cong 
\mathrm{Hom}_{\mathbf{Expr}}(E, K(G)).
\]
\end{theorem}

\section{Commutative Diagram}

\begin{center}
\begin{tikzcd}
E \arrow[r,"f"] \arrow[d,"\eta_E"'] & E' \arrow[d,"\eta_{E'}"] \\
\Pi(E) \arrow[r,"\Pi(f)"'] & \Pi(E')
\end{tikzcd}
\end{center}

\chapter{Verification Layer}

\section{Bisimulation}

\begin{definition}[Kernel Bisimulation]
Define $R$ by $(g_1,g_2) \in R$ iff $K(g_1) = K(g_2)$.
\end{definition}

\begin{theorem}
$R$ is a bisimulation on $G_M$.
\end{theorem}

\chapter{References}

\begin{thebibliography}{99}

\bibitem{LeeSmooth}
J.~M.~Lee, \emph{Introduction to Smooth Manifolds}, Springer, 2012.

\bibitem{Spivak}
M.~Spivak, \emph{A Comprehensive Introduction to Differential Geometry}, 
Publish or Perish, 1979.

\bibitem{MacLane}
S.~Mac Lane, \emph{Categories for the Working Mathematician}, Springer, 1998.

\bibitem{Riehl}
E.~Riehl, \emph{Category Theory in Context}, Dover, 2016.

\bibitem{AmbroseSinger}
W.~Ambrose and I.~M.~Singer, ``A Theorem on Holonomy,'' 
\emph{Trans. Amer. Math. Soc.}, 1953.

\bibitem{Hatcher}
A.~Hatcher, \emph{Algebraic Topology}, Cambridge University Press, 2002.

\bibitem{May}
J.~P.~May, \emph{A Concise Course in Algebraic Topology}, Chicago, 1999.

\bibitem{Lamport}
L.~Lamport, \emph{Specifying Systems}, Addison-Wesley, 2002.

\end{thebibliography}

\end{document}
